% !TEX program = lualatex
\documentclass[12pt, openany]{book}

\usepackage{fontspec}
\usepackage{geometry}
\geometry{
  paperwidth=6in,
  paperheight=9in,
  top=1.1in,
  bottom=1.1in,
  inner=0.9in,
  outer=0.85in,
}

\usepackage{microtype}
\emergencystretch=2em
\usepackage{setspace}
\usepackage{emptypage}
\usepackage{titlesec}
\usepackage{booktabs}
\usepackage{longtable}
\usepackage{array}
\usepackage{xcolor}
\usepackage{fancyhdr}
\usepackage{hyperref}
\usepackage{tikz}
\usetikzlibrary{positioning, arrows.meta, calc}

\definecolor{nodecolor}{gray}{0.92}
\definecolor{edgecolor}{gray}{0.45}

\hypersetup{
  colorlinks=false,
  pdfborder={0 0 0},
  pdftitle={Collected Works},
  pdfauthor={Flyxion},
}

\setmainfont{texgyrepagella-regular.otf}[
  BoldFont       = texgyrepagella-bold.otf,
  ItalicFont     = texgyrepagella-italic.otf,
  BoldItalicFont = texgyrepagella-bolditalic.otf,
]

\hyphenation{dis-tin-guish-a-bil-i-ty ad-mis-si-bil-i-ty phys-i-form-er sphere-pop}

\setstretch{1.35}

% Chapter headings — large scshape, no number, wide top margin
\titleformat{\chapter}[block]
  {\normalfont\large\scshape}
  {}
  {0pt}
  {}
\titlespacing*{\chapter}{0pt}{1.8in}{0.55in}

% Section headings — small scshape
\titleformat{\section}[block]
  {\normalfont\normalsize\scshape}
  {}
  {0pt}
  {}
\titlespacing*{\section}{0pt}{1.6\baselineskip}{0.5\baselineskip}

% Subsection headings — run-in italic
\titleformat{\subsection}[runin]
  {\normalfont\normalsize\itshape}
  {}
  {0pt}
  {}[.\enspace]
\titlespacing*{\subsection}{0pt}{0.8\baselineskip}{0pt}

% Running headers
\pagestyle{fancy}
\fancyhf{}
\fancyhead[LE]{\small\scshape flyxion}
\fancyhead[RO]{\small\scshape\leftmark}
\fancyfoot[C]{\small\thepage}
\renewcommand{\headrulewidth}{0pt}
\renewcommand{\chaptermark}[1]{\markboth{#1}{}}

% Plain style for chapter opening pages
\fancypagestyle{plain}{%
  \fancyhf{}
  \fancyfoot[C]{\small\thepage}
  \renewcommand{\headrulewidth}{0pt}
}

% TOC formatting
\usepackage{tocloft}
\renewcommand{\cftchapfont}{\normalfont\scshape}
\renewcommand{\cftchappagefont}{\normalfont}
\renewcommand{\cftchapdotsep}{\cftdotsep}
\setlength{\cftbeforechapskip}{0.5\baselineskip}
\renewcommand{\contentsname}{\normalfont\scshape\large contents}

\begin{document}

% ---------------------------------------------------------------
% HALF TITLE
% ---------------------------------------------------------------
\pagestyle{empty}
\vspace*{2.8in}
{\centering
\fontsize{11}{14}\selectfont\textsc{collected works}
\par}
\cleardoublepage

% ---------------------------------------------------------------
% TITLE PAGE
% ---------------------------------------------------------------
\thispagestyle{empty}
\vspace*{1.6in}
{\centering
\fontsize{22}{28}\selectfont\textit{Collected Works}

\vspace{0.55in}

\fontsize{14}{18}\selectfont\textsc{Flyxion}
\par}
\cleardoublepage

% ---------------------------------------------------------------
% TABLE OF CONTENTS
% ---------------------------------------------------------------
\pagestyle{fancy}
\tableofcontents
\cleardoublepage

% ---------------------------------------------------------------
% SECTION I — THE QUESTION (opening statement)
% ---------------------------------------------------------------
\chapter[The Question]{Section I \\ The Question}
\markboth{section i: the question}{}

\noindent Every structure inherits a problem.

\medskip

\noindent The conditions that produced it eventually disappear.
Authors die. Languages drift. Software becomes obsolete. Institutions decay.
Archives fragment. Memories fade. Environments change. The future withdraws
the assumptions on which the present was built.

\medskip

\noindent Yet some structures persist.

\medskip

\noindent This collection is an investigation of that persistence.

\bigskip

\noindent Its subjects range across mathematics, cognition, language, artificial
intelligence, software, archives, fiction, institutions, and civilization.
The common question is not what these things are, but how they continue.
What remains when the conditions that originally produced a structure are no
longer available? What forms of repair become possible after degradation?
What distinctions survive compression, translation, neglect, or loss?
Which futures remain admissible after pathways have collapsed?

\bigskip

\noindent The essays, systems, visualizations, stories, tools, and experiments
collected here should therefore be understood as observations of a single
phenomenon viewed through different instruments.

\medskip

\noindent Their central concern is the persistence of distinctions under the
withdrawal of the conditions that made those distinctions possible.

\bigskip\bigskip

\noindent The question remains open. That is why the work continues.

\cleardoublepage

% ---------------------------------------------------------------
% SECTION II — THE DISCOVERY
% ---------------------------------------------------------------
\chapter[The Discovery]{Section II \\ The Discovery}
\markboth{section ii: the discovery}{}

\noindent\textit{How the question emerged.}

\bigskip

\noindent The research did not begin with a question about survival.

It began closer to questions about information, representation, and
computation. Early work investigated what information is present in a
system, how states are encoded, how machines process symbols, how meaning
is represented in language and neural networks.

Those were reasonable starting points. They are the standard starting points.

But they produced recurring failures.

\section{The Failures}

Representation did not explain persistence. A system can possess a complete
representation of something and still lose it. The representation is not the
thing. Its accuracy at one moment does not guarantee recoverability at the next.

Prediction did not explain repair. A system optimized to predict the future
has no obvious mechanism for reconstructing a damaged past. Prediction and
repair are different operations, and the difference matters.

Optimization did not explain continuation. A system maximizing an objective
function will follow the gradient. It has no principled account of what to
do when the landscape becomes inadmissible — when the futures required by
the objective are no longer reachable.

State descriptions did not explain historical dependence. The same state can
be reached by different histories. Those histories are not equivalent. The
meaning of a current state often depends entirely on how it was reached.

Information did not explain why some distinctions survive while others
disappear. Information theory counts bits. It does not explain why some
boundaries remain stable under noise, compression, and transformation while
others collapse.

Each failure suggested the same correction: the problem was not with the
answers but with the level of description.

\section{The Phases}

The vocabulary shifted in phases.

In the first phase, the central question was representational: what
information is present? This produced work on encoding, semantics, machine
learning, and computational state.

In the second phase, the question became ontological: what distinctions make
information possible? Entities ceased to be primitive. Distinctions became
primitive. A word, a species, a government, a theorem — each exists because
some distinctions are maintained while others are ignored.

In the third phase, the question became restorative: how are distinctions
reconstructed after loss? Repair emerged as a first-class operation rather
than a secondary process. The repaired structure is not identical to the
original. Repair is constructive.

In the fourth phase, the question became dynamic: what allows structures to
persist through degradation? Continuation replaced conservation as the
central concept. Identity was relocated from substance to trajectory.

In the fifth phase, the question became geometric: which futures remain
reachable? Admissibility structures replaced state descriptions as the
primary objects of study. The most important property of a system became not
its present configuration but the shape of the futures it could still access.

In the sixth phase, these threads converged.

\section{The Convergence}

The convergence was not planned. It was observed.

Two independent reconstructions — one tracing the history of the projects,
one classifying the research programs — arrived at essentially the same
sentence:

\bigskip
\begin{center}
\textit{How do structures survive the loss of the conditions\\
that originally produced them?}
\end{center}
\bigskip

That convergence is itself evidence. A question that emerges independently
from multiple directions of inquiry is more likely to be the real question
than one stipulated at the outset.

The programs, instruments, experimental habitats, and accumulated corpus
described in the sections that follow are best understood as responses to
that question — each from a different direction, at a different scale,
using different instruments.

None of them fully answers it.

That is why the work continues.

\cleardoublepage

% ---------------------------------------------------------------
% SECTION III — THE PROGRAMS
% ---------------------------------------------------------------
\chapter[The Programs]{Section III \\ The Programs}
\markboth{section iii: the programs}{}

\noindent\textit{Fifteen research programs constituting the corpus.}

\bigskip

\noindent These programs did not arrive simultaneously. They accumulated
through the sequence of explanatory failures described in the preceding
section. Each one opened when the previous vocabulary proved insufficient.
Together they constitute a distributed investigation of a single phenomenon
approached from fifteen directions.

\section{I. The Ontological Critique Program}

The Ontological Critique Program investigates the creation, maintenance,
degradation, repair, and disappearance of distinctions.

It begins from the observation that every object, category, concept,
organism, institution, language, and scientific theory exists only because
some distinctions are maintained while others are ignored. Rather than
taking entities as primitive, the program treats distinctions as primitive.

A mountain is a distinction between elevation regimes. A species is a
distinction maintained through reproductive and ecological processes. A word
is a distinction stabilized through repeated social use. A government is a
distinction maintained by administrative and legal mechanisms. A scientific
theory is a distinction between admissible and inadmissible explanations.

The central question is not what things are, but how distinctions become
stable enough to appear as things.

Memory becomes the preservation of distinctions through time. Repair becomes
the reconstruction of lost distinctions. Information becomes structured
distinction. Entropy becomes the erosion or equivalence of distinctions. The
long-term objective is a general theory of distinguishability applying
equally to physics, cognition, institutions, biology, language, and
computation.

\section{II. The Admissibility Program}

The Admissibility Program investigates the geometry of possibility.

Most theories begin by asking what exists. This program instead asks what
remains possible. A system is characterized not merely by its current state
but by the future states that remain reachable from it. These reachable
futures form an admissibility structure.

A civilization possesses admissibility. A language possesses admissibility.
A brain possesses admissibility. A software project possesses admissibility.
The destruction of admissibility often precedes visible collapse. A city can
appear prosperous while future options disappear. A scientific field can
appear productive while alternative research trajectories become inaccessible.

The program develops admissibility manifolds, reachability volumes,
admissibility distortion metrics, admissibility equivalence relations, and
boundary proximity measures. Its central claim is that the most important
property of a system is often not its present configuration but the
structure of futures it can still access.

\section{III. The Repair Theory Program}

Repair Theory studies how systems recover after damage.

Traditional science often focuses on equilibrium, optimization, or
prediction. Repair Theory focuses on restoration, regeneration, adaptation,
and continuation. A repaired system is rarely identical to its previous
state. Instead, repair creates new pathways through a damaged landscape.

The program investigates biological repair, institutional repair, cognitive
repair, linguistic repair, computational repair, and civilizational repair.
A recurring result is that repair is constructive rather than conservative.
The repaired bridge differs from the original bridge. The repaired language
differs from the original language. The repaired memory differs from the
original memory. Repair therefore becomes a source of novelty rather than
merely a return to equilibrium.

Many later works treat memory itself as a repair process and intelligence as
a special form of repair operating over prediction failures.

\section{IV. The Continuation Program}

The Continuation Program asks what it means for something to remain itself
through change.

Traditional identity theories often emphasize static properties. Continuation
theory instead focuses on trajectories. A river remains a river despite
exchanging water. A city remains a city despite replacing buildings. A
language remains a language despite changing vocabulary. A person remains a
person despite biological turnover.

Identity is relocated from substance to continuation. The program studies
path preservation, reachability maintenance, structural persistence,
historical continuity, and dynamic identity. Its central intuition is that
persistence is not conservation of matter but conservation of continuation.

\section{V. The Memory Program}

The Memory Program begins with a rejection of representational primacy.

Memory is not fundamentally an image, symbol, representation, or stored
object. Memory is a recoverability structure. A system remembers when
previous distinctions can be reconstructed. This view applies across
biological memory, archives, scientific records, databases, cultural
traditions, and language.

Memory therefore becomes a geometric property of histories rather than a
collection of stored representations. The program investigates event logs,
recoverability, ecphory, compression, reconstruction, and historical
persistence. Its most important claim is that memory is fundamentally about
continuation rather than storage.

\section{VI. The Document Ecology Program}

The Document Ecology Program studies documents as living persistence
structures.

Most publishing systems treat documents as static containers. This program
treats documents as organisms moving through technological environments. A
document must survive format migration, compression, copying, institutional
neglect, software obsolescence, hardware failure, and linguistic drift.

Icepick, Marqants, readability-first publishing, snapshot systems, and
archival workflows all belong to this program. The central question becomes:
how can knowledge survive the disappearance of the systems that originally
encoded it? This transforms publishing from presentation into a problem of
long-term persistence.

\section{VII. The Coordination Geometry Program}

The Coordination Geometry Program studies how many agents become coherent.

Rather than focusing on individual decision-makers, it investigates
collective structure across markets, governments, open-source communities,
scientific fields, neural populations, and social movements. The program
examines synchronization, information transport, coordination failures,
institutional coherence, and collective repair.

Its central concern is understanding how distributed systems generate unified
behavior without centralized control.

\section{VIII. The Preference Fields Program}

The Preference Fields Program attempts to geometrize value.

Instead of treating preferences as isolated decisions, it models them as
fields defined over possibility spaces. An agent does not merely choose among
options. It inhabits a landscape where some futures possess higher preference
density than others. The program studies preference manifolds, value
gradients, collective preference structures, dynamic preference evolution,
and preference transport.

The long-term objective is a mathematical account of decision-making
compatible with the admissibility framework.

\section{IX. The Computation-as-History Program}

This program arises from Spherepop, history-based programming, and related
systems.

Traditional computation emphasizes state. This program emphasizes history.
The meaning of a computation depends not merely on where it arrives but on
how it arrived there. Key concepts include event histories, refusal,
collapse, operational trajectories, and historical semantics. Programs become
evolving histories rather than static transformations.

\section{X. The Procedural Attention Program}

This program investigates how experiences are designed.

Cinema, games, music, textbooks, proofs, and interfaces all guide observers
through structured possibility spaces. Meaning emerges from traversal. The
program studies attention routing, gaze control, narrative geometry,
procedural reward, and traversal systems.

Its central claim is that designed experiences are fundamentally systems for
generating trajectories through admissibility structures.

\section{XI. The Fictional Habitat Program}

The Fictional Habitat Program includes City of Brutes, Notes from the
Collapsewell, Worlds of If, the glossary-cycle literature, and the
Instrumentarium. These works are not illustrations of theory. They are
experimental environments.

Rather than arguing for concepts, they instantiate them. A degrading
orthography becomes a model of persistence under corruption. A fictional
ecology becomes a model of distinction formation. An ergodic text becomes a
repair problem. The reader becomes a participant in the phenomenon under
investigation.

This program therefore serves as an experimental laboratory for theoretical
concepts — a place where the survival question is not argued but enacted.

\section{XII. The Visualization Program}

The Visualization Program attempts to make invisible structures observable.

Examples include Audioscope, the Lava Thought Visualizer, the Topological
Loop Monitor, reachability landscapes, educational comics, and semantic
terrain displays. The objective is not artistic decoration. The objective is
witness extraction.

Visualization functions as instrumentation. It allows structures that are
difficult to perceive verbally or mathematically to become directly
inspectable.

\section{XIII. The Civilization and Futures Program}

This program investigates long-term persistence at civilizational scales.

Its subjects include infrastructure, institutions, energy systems, governance,
technological progress, fiscal reachability, and state capacity. The central
question is which civilizations preserve admissible futures and which consume
them.

The program interprets collapse not as sudden failure but as progressive
contraction of reachable futures. Civilizations are evaluated according to
their ability to maintain repair capacity, continuation pathways, and future
option space.

\section{XIV. The Generative Evolution Program}

The Generative Evolution Program investigates open-ended novelty generation
and design-space exploration.

Its principal projects include Blastoids, pic-breeder style evolutionary
systems, and morphological search procedures. The program studies
evolutionary aesthetics, fitness landscape geometry, and the conditions
under which open-ended novelty is possible rather than convergent
optimization.

It is structurally distinct from the other programs in that its primary
concern is not the survival of existing structures but the generation of
new ones — though the two questions turn out to be deeply connected. A
system capable of genuine novelty generation is also a system capable of
finding repair pathways through previously unoccupied regions of a damaged
landscape.

\section{XV. The RSVP Program}

The RSVP Program — Relativistic Scalar-Vector Plenum — is the primary
physics program and is structurally different from the fourteen programs
above.

Most of the other programs are epistemological, ontological, computational,
informational, institutional, or methodological. RSVP is the only major
program attempting a direct account of physical reality itself.

RSVP proposes an alternative cosmological framework in which scalar entropy
fields and vector baryon flows replace metric expansion as the organizing
principle of cosmic evolution. The framework addresses structure formation,
the cosmic microwave background, entropic smoothing mechanisms, and
alternatives to inflationary cosmology.

Within the RSVP program the collaboration has developed scalar entropy
fields, vector baryon flows, CTOV equations, entropic gravity analogues,
derived algebraic geometric formulations using shifted symplectic geometry
and BV-BRST formalisms, consciousness metrics within the field-theoretic
framework, and experimental predictions distinguishing RSVP from standard
cosmology.

The program asks: what is capacity? What is transport? What is entropy? How
does structure emerge physically? What governs cosmological evolution? These
are questions at a different level than the other programs — not about the
geometry of cognition or the persistence of institutions but about the
geometry of the universe itself.

RSVP is included here as the fifteenth program because the same
distinguishability question that organizes all the others reappears at the
cosmological scale. Cosmic structure is itself a distinction maintained
against entropic erosion. The universe too is a system that must solve the
survival problem.

\bigskip\bigskip

\noindent\textit{Together these fifteen programs constitute a distributed
investigation. Each addresses the survival question from a different
direction, at a different scale, using different instruments. The
instruments themselves are catalogued in the section that follows.}

\cleardoublepage

% ---------------------------------------------------------------
% SECTION IIIA — PROGRAM RELATIONS
% ---------------------------------------------------------------
\chapter[Program Relations]{Section IIIA \\ Program Relations}
\markboth{section iiia: program relations}{}

\noindent\textit{How the fifteen programs depend on and generate one another.}

\bigskip

\noindent The fifteen programs described in Section III are not independent.
They form a dependency structure in which some programs supply the
foundational vocabulary on which others are built, and in which certain
programs generate phenomena that other programs are then designed to study.

Understanding that structure is more useful than reading the programs as a
flat list.

\section{The Principal Chain}

One sequence runs through nearly every major work in the corpus. It is not
a logical derivation. It is a historical and conceptual dependency: each
link in the chain presupposes the previous one and motivates the next.

\bigskip

\begin{center}
\begin{tikzpicture}[
  node distance=0.52in,
  every node/.style={
    rectangle,
    rounded corners=3pt,
    fill=nodecolor,
    draw=edgecolor,
    text width=1.55in,
    align=center,
    font=\small\itshape,
    inner sep=6pt,
  },
  arrow/.style={
    -Stealth,
    thick,
    color=edgecolor,
  }
]
\node (dist) {Distinction};
\node (mem)  [below=of dist] {Memory};
\node (rep)  [below=of mem]  {Repair};
\node (cont) [below=of rep]  {Continuation};
\node (adm)  [below=of cont] {Admissibility};
\node (surv) [below=of adm]  {Survival};
\draw[arrow] (dist) -- (mem);
\draw[arrow] (mem)  -- (rep);
\draw[arrow] (rep)  -- (cont);
\draw[arrow] (cont) -- (adm);
\draw[arrow] (adm)  -- (surv);
\end{tikzpicture}
\end{center}

\bigskip

The chain reads as follows.

Distinction is the most primitive operation. Before anything can be
remembered, repaired, continued, or evaluated for admissibility, it must
first be distinguished. The Ontological Critique Program supplies this
foundation.

Memory is the persistence of distinctions through time. A system that
cannot maintain distinctions across intervals cannot accumulate the
history required for repair. The Memory Program builds directly on the
distinction framework and is unintelligible without it.

Repair is what a system does when memory is incomplete or damaged. It
reconstructs distinctions from available evidence. Repair presupposes
memory because repair requires knowing what has been lost. The Repair
Theory Program therefore sits above the Memory Program in the dependency
order.

Continuation is the outcome of successful repair extended across time. A
system that repairs isolated distinctions but cannot thread those repairs
into an ongoing trajectory does not continue; it merely restarts. The
Continuation Program takes repair as its input and asks under what
conditions repair produces trajectory rather than mere restoration.

Admissibility is the geometry of the futures that remain reachable after
a history of distinctions, memories, repairs, and continuations. The
Admissibility Program asks which of those futures are still accessible
and which have been foreclosed. It presupposes continuation because
admissibility is a property of trajectories, not states.

Survival is the terminal question to which the chain leads. A structure
survives when its distinctions persist, its memories remain recoverable,
its repairs produce continuation, and its admissible futures remain
non-empty. The chain is the formal anatomy of survival.

\section{Lateral Dependencies}

Several programs depend on the principal chain without being part of it.

The \textit{RSVP Program} sits beneath the entire chain as a physical
substrate. It provides the field-theoretic account of what distinction,
entropy, and reachability mean at cosmological scale. The chain describes
what must happen for a structure to survive; RSVP describes the physical
medium in which that survival is negotiated.

The \textit{Computation-as-History Program} is a formalization of the
Memory link in the chain. It treats computation itself as a
history-dependent process and inherits the full dependency structure from
Memory upward.

The \textit{Coordination Geometry Program} depends on Continuation and
Admissibility. A collective can only coordinate if its members share
admissible continuation pathways. Coordination failure is often
admissibility failure at the collective scale.

The \textit{Preference Fields Program} depends on Admissibility.
Preferences are meaningful only within an admissibility structure: an
agent cannot genuinely prefer a future that is not reachable. Preference
geometry is therefore a specialization of admissibility geometry.

The \textit{Civilization and Futures Program} applies the full chain at
civilizational scale. A civilization is a structure that must solve the
survival problem at the level of institutions, infrastructure, governance,
and collective memory.

\section{Methodological Programs}

Several programs are not positions in the dependency chain but methods
for studying it.

The \textit{Procedural Attention Program} is a method for studying how
observers traverse admissibility structures. Cinema, games, and interfaces
are all instruments for generating trajectories through designed possibility
spaces.

The \textit{Visualization Program} is a method for making the chain
inspectable. Admissibility manifolds, reachability landscapes, and
distinction field dynamics are all difficult to perceive without
instrumentation. The visualization projects provide that instrumentation.

The \textit{Document Ecology Program} is a method for preserving the
chain's outputs. Papers, monographs, and systems are themselves structures
that must solve the survival problem. The document ecology program applies
the chain to its own products.

The \textit{Generative Evolution Program} studies what happens when repair
is replaced by novelty generation — when a damaged landscape is navigated
not by restoring what was lost but by finding what has never existed. It
sits outside the principal chain but connects to Admissibility: novelty
generation is the exploration of admissibility space rather than the
recovery of prior positions within it.

\section{The Fictional Habitats as Experimental Closure}

The \textit{Fictional Habitat Program} occupies an unusual structural
position. It does not depend on the chain in the way the analytical
programs do. Instead, it provides experimental closure.

The analytical programs produce claims about what happens when distinctions
degrade, memories fail, repair occurs, continuation succeeds or fails, and
admissibility contracts. Those claims are formally derived but
phenomenologically unverified. The fictional habitats create conditions
under which those phenomena can be directly experienced by a reader.

The fictional habitats therefore close the loop between the formal chain
and the phenomena the chain describes. They are neither foundational nor
methodological but experimental: environments in which the survival
question is enacted rather than argued.

\bigskip

\noindent\textit{The chain — Distinction, Memory, Repair, Continuation,
Admissibility — is the conceptual skeleton of the corpus. Everything else
either instantiates a link in that chain, studies it from outside, or
provides the physical or experimental substrate in which it operates.}

\cleardoublepage

% ---------------------------------------------------------------
% SECTION IV — THE INSTRUMENTS
% ---------------------------------------------------------------
\chapter[The Instruments]{Section IV \\ The Instruments}
\markboth{section iv: the instruments}{}

\noindent\textit{A catalog of the observational and constructive devices.}

\bigskip

\noindent The corpus is not merely a collection of theories. It is a
collection of instruments. Many of the named systems function less like
finished accounts and more like microscopes — each aimed at a different
aspect of persistence, continuation, and the survival of structure.

The distinction between a theory and an instrument matters here. A theory
makes claims. An instrument makes phenomena inspectable. Several of the
systems listed below do both, but their primary function is instrumental:
they allow structures that would otherwise be invisible to become directly
observable, traversable, or manipulable.

The instruments are organized by type.

\section{Physical Instruments}

These instruments are directed at physical reality. They are the primary
tools of the RSVP Program.

\subsection{RSVP} Relativistic Scalar-Vector Plenum. The principal
cosmological framework. Models cosmic evolution through scalar entropy fields
and vector baryon flows rather than metric expansion. Generates CTOV
equations governing the coupled dynamics of capacity, transport, orientation,
and volume. Functions as the physical substrate into which the admissibility
and distinguishability frameworks embed at cosmological scale.

\subsection{CTOV equations} Capacity, Transport, Orientation, Volume. The
governing equations of the RSVP framework. Analogous in structure to
Maxwell's equations but operating over entropic rather than electromagnetic
fields.

\subsection{Admissibility manifolds} Formal geometric objects encoding which
future states remain reachable from a given configuration. Defined over
arbitrary state spaces. Apply equally in physical, cognitive, institutional,
and computational contexts.

\subsection{Reachability volumes} The measure of future option space
available to a system at a given moment. Contraction of reachability volume
is the formal signature of approaching collapse — in physical systems,
institutions, and computational processes alike.

\section{Distinction Instruments}

These instruments are directed at the formation, maintenance, and
degradation of distinctions.

\subsection{Distinguishability geometry} The formal framework treating
distinguishability as a geometric property of state spaces. Two states are
distinguishable to the degree that they can be separated by an observer
operating within given constraints. Provides a common mathematical language
connecting ontology, computation, and physics.

\subsection{Ontological deficit} A measure of the degree to which a system
has lost its capacity to maintain distinctions that were previously
maintainable. Analogous to information loss but defined in terms of
boundary stability rather than bit count.

\subsection{Collapse quotient} The ratio of distinctions lost to distinctions
maintained across a transformation. Governs the information geometry of
compression, projection, and degradation.

\subsection{Distinction field} A field defined over a possibility space
assigning to each region a local distinguishability coefficient. Dense
regions support fine distinctions. Sparse regions support only coarse ones.

\section{Repair Instruments}

These instruments are directed at the reconstruction of damaged or degraded
structures.

\subsection{Repair algebra} A formal algebraic system encoding the operations
available to a system attempting to reconstruct lost distinctions. The
algebra is non-commutative: the order of repair operations matters.
Reconstruction via pathway A produces a different result than reconstruction
via pathway B even when both begin and end at nominally the same states.

\subsection{Persistence operators} Operators acting on state histories that
extract the components of a structure most likely to survive degradation.
Persistence operators identify which distinctions are load-bearing and which
are redundant.

\subsection{Continuation function} A function mapping damaged states to the
set of continuation pathways still available. The continuation function is
the primary object studied in the Repair Theory and Continuation programs.

\subsection{Ecphory} The process by which a partial cue reconstructs a
fuller memory or distinction. Treated in the Memory Program as the primary
mechanism of recoverability. Distinguished from retrieval, which presupposes
intact storage.

\section{Memory and Archive Instruments}

These instruments are directed at the persistence of histories.

\subsection{MEM\textbar8} A history-centered computational architecture
emphasizing persistent memory and structural continuity. Treats computation
as fundamentally a process of maintaining and extending event logs rather
than transforming isolated states.

\subsection{Event log} The primary data structure of history-based
computation. Records not just outcomes but the trajectories by which
outcomes were reached. Provides the substrate for ecphory, repair, and
continuation.

\subsection{Recoverability structure} A formal characterization of which
past distinctions remain reconstructible from a given present state and set
of available cues. The central object of the Memory Program.

\subsection{Marqants} A publishing markup system designed for long-term
document survival. Prioritizes recoverability over presentational richness.
Encodes semantic structure in forms that survive format migration, software
obsolescence, and institutional neglect.

\section{Computational Instruments}

These instruments are directed at computation conceived as a historical
process.

\subsection{Spherepop} An irreversible event calculus and computational
substrate. Treats programs as evolving event histories rather than state
machines. Implements refusal as a first-class operation: a computation can
decline to proceed, and that refusal is itself an event with semantic content.

\subsection{cprsh} A shell environment implementing constraint-first
computation. Commands operate by declaring constraints on reachable states
rather than specifying transformations of current states.

\subsection{History machines} Computational architectures in which the full
event history of a process is a first-class data structure, accessible and
manipulable at runtime.

\subsection{Refusal calculus} A formal account of the semantics of
computational refusal. Specifies the conditions under which a process
legitimately declines to proceed and what information is conveyed by that
declination.

\subsection{Icepick} A codebase navigation and snapshot system. Treats a
software repository not as a collection of files but as a traversable
historical object. Preserves not just current states but the trajectories
through which current states were reached.

\section{AI and Cognitive Instruments}

These instruments are directed at artificial and natural cognition.

\subsection{HYDRA} A recursive computational organization framework for
distributed cognitive architectures. Addresses how large systems maintain
coherent behavior across heterogeneous subsystems with different memory
structures and update rates.

\subsection{CLIO} A cognitive architecture emphasizing longitudinal memory
and the accumulation of historical context across extended interactions.
Designed to support reasoning processes that depend on trajectory rather
than momentary state.

\subsection{TARTAN} A recursive tiling and spatial organization framework.
Implements Gray-code tilings, L-system derived geometry, and
entropy-preserving tiling dynamics. Functions as an organizational substrate
for spatial reasoning and semantic terrain display.

\subsection{Admissibility distortion metric} A measure of how severely a
cognitive or computational system's future option space has been deformed
by accumulated commitments, compressed representations, or pathway closures.

\section{Visualization Instruments}

These instruments make invisible structures directly inspectable.

\subsection{Audioscope} An audio visualization system using CRT and phosphor
aesthetics to render sound as spatial trajectory. The visual form embeds
the temporal structure of the signal rather than merely displaying amplitude
over time.

\subsection{Lava Thought Visualizer} A visualization system rendering
semantic terrain as dynamic topographic landscape. Conceptual proximity
appears as spatial proximity. Conceptual drift appears as landscape
deformation.

\subsection{Topological Loop Monitor} A visualization system tracking the
formation, persistence, and destruction of topological features in dynamical
systems. Renders persistence diagrams in real time.

\subsection{Reachability landscapes} Visualizations of admissibility
structures as navigable terrain. High ground represents accessible futures.
Cliffs represent catastrophic admissibility boundaries. Valley floors
represent attractors.

\subsection{Semantic terrain displays} Visualizations of semantic spaces
as navigable geographic environments. Used in the Procedural Attention
Program to study how readers and viewers traverse conceptual landscapes.

\subsection{Educational comics} A recurring visual format combining
diagrammatic precision with narrative structure. Used to make formal
theoretical content accessible without sacrificing accuracy. The retro
infographic aesthetic encodes a deliberate epistemic stance: knowledge
should be transmissible across technological generations.

\section{Publishing and Preservation Instruments}

These instruments are directed at the long-term survival of documents
and knowledge.

\subsection{Readability-first publishing} A publishing philosophy
prioritizing legibility across future reading environments over
presentational sophistication in the present. Treats typographic simplicity
as a survival strategy.

\subsection{Randomized compilation artifacts} Documents produced through
randomized \LaTeX{} compilation procedures that introduce controlled
variation into layout, spacing, and typographic decisions. Serve as
experiments in the aesthetics of document survival.

\subsection{Self-archiving systems} Workflows encoding provenance,
versioning, and reconstruction instructions directly into documents.
A document that carries its own archival metadata is more likely to survive
institutional discontinuity than one that depends on external cataloguing
systems.

\section{Evolutionary and Generative Instruments}

These instruments are directed at novelty generation and design-space
exploration.

\subsection{Blastoids} An evolutionary image generation system implementing
pic-breeder style selection. Users navigate a fitness landscape through
aesthetic preference rather than explicit objective specification.
Demonstrates that open-ended novelty is possible in bounded search spaces
when selection pressure is human rather than algorithmic.

\subsection{Morphological search} Search procedures operating over spaces
of structural forms rather than parameter spaces. Explores regions of
design space inaccessible to gradient-based methods.

\bigskip\bigskip

\noindent The table below provides a compact overview of the instrument
catalog organized by program.

\bigskip

\begin{longtable}{@{} p{0.28\textwidth} p{0.68\textwidth} @{}}
\toprule
\textsc{Program} & \textsc{Primary Instruments} \\
\midrule
Physical & RSVP, CTOV, Admissibility manifolds, Reachability volumes \\[4pt]
Distinction & Distinguishability geometry, Ontological deficit, Collapse quotient \\[4pt]
Repair & Repair algebra, Persistence operators, Continuation function \\[4pt]
Memory & MEM\textbar8, Ecphory, Recoverability structure, Event log \\[4pt]
Computation & Spherepop, cprsh, History machines, Refusal calculus, Icepick \\[4pt]
AI / Cognitive & HYDRA, CLIO, TARTAN, Admissibility distortion metric \\[4pt]
Visualization & Audioscope, Lava Thought Visualizer, Topological Loop Monitor \\[4pt]
Publishing & Marqants, Readability-first publishing, Self-archiving systems \\[4pt]
Evolutionary & Blastoids, Morphological search \\
\bottomrule
\end{longtable}

\bigskip

\noindent\textit{The instruments in each row are not merely associated with
the program listed. They are the means by which that program makes its
phenomena inspectable. Without the instruments, the programs would remain
purely verbal. Without the programs, the instruments would have no
orientation.}

\cleardoublepage

% ---------------------------------------------------------------
% SECTION V — EXPERIMENTAL HABITATS
% ---------------------------------------------------------------
\chapter[Experimental Habitats]{Section V \\ Experimental Habitats}
\markboth{section v: experimental habitats}{}

\noindent\textit{Fiction, narrative environments, bestiaries, ergodic
texts, and comics as methodology.}

\bigskip

\noindent Once the preceding sections are in place — the question, the
history of its discovery, the fifteen programs, the instrument catalog —
the works described in this section become legible in a way they would
not be otherwise.

They are not side projects. They are not illustrations of the theoretical
work. They are not relaxations from it.

They are experiments conducted in a different medium.

The theoretical programs argue that some structures survive while others
do not. The fictional habitats enact survival and failure directly. A
reader navigating a degrading orthography is not reading about persistence
under corruption. The reader is experiencing it. The medium is the
phenomenon.

This methodological distinction matters because it produces evidence that
the formal programs cannot produce on their own. A mathematical account of
distinction degradation can specify what is lost and at what rate. It
cannot specify what the loss feels like from the inside of the degrading
structure. The fiction can.

The habitats described below are therefore best understood as observational
environments — constructed spaces in which the phenomena of persistence,
repair, continuation, admissibility, and survival are allowed to occur and
can be directly inspected by a participant observer.

\section{Primary Narrative Environments}

\subsection{City of Brutes} A screenplay. The primary setting is a city
organized around competing libraries: the Library of Revisions and the
Library of Contradictions. The central tension is between institutions that
preserve distinctions and institutions that dissolve them. Characters
navigate an urban environment in which the maintenance of meaning requires
continuous active effort against ambient institutional entropy.

The city functions as a model of civilizational admissibility. The
screenplay traces what happens when repair capacity is exhausted before
the structures requiring repair have finished degrading. It does not
argue that this is how cities fail. It constructs a space in which a reader
can observe the failure from within.

\subsection{Notes from the Collapsewell} An ergodic novel. The text itself
undergoes controlled degradation across the course of the reading. Spelling
systems drift. Vocabulary shifts. Grammatical structures erode and
reconstitute. The reader's experience of the novel at page two hundred
is structurally different from the experience at page twenty, not because
the narrative has advanced but because the medium has changed.

This is a direct experimental test of a claim made in the Memory Program:
that meaning can survive the degradation of the vehicle that carries it.
The novel answers the question empirically rather than argumentatively.
Either the reader continues to construct coherent meaning from a
deteriorating text or the reader does not. The outcome is data.

\subsection{Worlds of If} A series of speculative documentary texts in the
style of alternative history magazines. Each issue investigates a trajectory
that was possible but did not occur — a reachable future that was not
reached. The series functions as applied admissibility theory. It asks
what conditions would have been sufficient to access a foreclosed future,
and what specific failure of repair, continuation, or distinction
maintenance caused that future to become inaccessible.

The documentary format is deliberate. Worlds of If does not present
alternative histories as fiction in the conventional sense. It presents
them as recovered evidence from trajectories that were abandoned. The
epistemological posture is archaeological rather than imaginative.

\section{Ecological and Taxonomic Environments}

\subsection{The Flyxion Instrumentarium} A bestiary of conceptual entities.
The Instrumentarium catalogs the creatures, instruments, and ecological
structures that inhabit the theoretical landscape of the research programs.

Entries include Witnesses, Repairers, Distinction-generating entities,
Admissibility monitors, Continuation organisms, and Preference field
inhabitants. Each entry follows the conventions of natural history writing:
habitat, behavior, diet, distinguishing characteristics, ecological role,
and threat status.

The bestiary format is not decorative. Natural history writing developed
precise conventions for describing entities whose behavior can be observed
but whose inner states are inaccessible. Those conventions translate
directly to the problem of describing formal theoretical entities whose
behavior can be derived but whose phenomenology is unknown.

The Instrumentarium sits at the intersection of philosophy, folklore,
systems theory, and scientific illustration. It is the only part of the
corpus that attempts to give the theoretical entities habitats — to
situate them in environments where they must compete, cooperate, reproduce,
and survive.

\subsection{The Distinction Ecology} A companion project to the
Instrumentarium. Where the Instrumentarium catalogs individual entities,
the Distinction Ecology describes the environments those entities inhabit.
It maps the terrain of conceptual space as a biological environment:
identifying niches, food webs, extinction pressures, and the conditions
under which new distinctions can emerge and stabilize.

\section{Glossary and Lexical Environments}

\subsection{Glossary-cycle literature} A series of texts organized around
evolving glossaries. Each work begins with a controlled lexicon and tracks
the semantic drift of that lexicon across the course of the text. Terms
acquire new meanings through use. Old meanings become inaccessible. The
text functions as a demonstration of the Continuation Program's central
claim: that identity is a property of trajectories rather than substances.

The glossary-cycle format makes semantic drift visible in a way that
ordinary narrative cannot. Because the terms are explicitly defined at the
outset, deviations from those definitions are legible as events rather than
noise.

\subsection{Intentional spelling degradation} A system of controlled
orthographic corruption developed for use in Notes from the Collapsewell
and related texts. The degradation is not random. It follows rules derived
from the distinction geometry framework: the boundaries most likely to
survive are those that are multiply reinforced, while singly reinforced
boundaries erode first.

The spelling degradation system is therefore a physical implementation of
the Ontological Critique Program's central model. It demonstrates that
even orthographic distinctions have a persistence structure and that some
spellings are more admissible than others under degradation pressure.

\section{Visual and Sequential Environments}

\subsection{Landscape comics} A recurring visual format combining
diagrammatic representation of theoretical structures with sequential
narrative. The retro infographic aesthetic is deliberate: it encodes a
commitment to transmissibility across technological generations. A comic
drawn in the style of a 1960s educational publication does not require
software to read.

Landscape comics have been produced for several of the theoretical
programs, including visualizations of admissibility landscapes, reachability
terrain, distinction field dynamics, and repair pathway geometry. Each
functions as a habitat in miniature: a bounded environment in which the
phenomenon under investigation can be observed at a scale legible to a
single reader.

\subsection{The four-vocabulary civilization narrative} A narrative framework
organizing fictional civilizational history around four classes of
participant: Inhabitants, Historians, Repairers, and Narrators. Each class
possesses a different relationship to the distinctions that constitute the
civilization.

Inhabitants maintain distinctions through practice without necessarily
theorizing them. Historians reconstruct distinctions that have partially
degraded. Repairers reconstruct distinctions that have fully collapsed.
Narrators maintain the meta-distinction between what the civilization was,
what it is, and what it might yet become.

The four-vocabulary framework is used in City of Brutes, in several of
the Worlds of If issues, and in theoretical discussions of civilizational
admissibility. It provides a structural vocabulary for discussing the
division of epistemic labor within a persistence-maintaining community.

\section{The Epistemological Status of the Habitats}

The fictional and creative works in this section occupy an unusual
epistemological position relative to the theoretical programs.

They are not evidence in the conventional scientific sense. They do not
provide data that could falsify a formal hypothesis. But they are not
merely illustrative either. They provide a kind of evidence that formal
argument cannot provide: phenomenological evidence.

A formal account of distinction degradation specifies the rate at which
boundaries erode and the conditions under which recovery is possible. It
does not specify what it is like to be inside a degrading distinction
structure — to be a reader whose interpretive tools are themselves
degrading as the text proceeds.

The habitats provide that specification. They are experiments in the
original sense: controlled arrangements of conditions under which a
phenomenon can be observed. The laboratory is the text. The apparatus
is the narrative structure. The observer is the reader.

What the reader discovers — or fails to discover — in the course of
inhabiting these environments constitutes genuine evidence about the
phenomena the theoretical programs are investigating. Evidence of a
kind that does not reduce to argument and cannot be produced by
argument alone.

\bigskip\bigskip

\noindent\textit{The experimental habitats are the part of the corpus
most likely to survive in forms unrecognizable to the research programs
that generated them. A reader in a future that has forgotten the
admissibility framework may still navigate the Collapsewell. The
phenomenon will persist even when the vocabulary for describing it
has been lost. That is, of course, the point.}

\cleardoublepage

% ---------------------------------------------------------------
% SECTION VI — THE CORPUS
% ---------------------------------------------------------------
\chapter[The Corpus]{Section VI \\ The Corpus}
\markboth{section vi: the corpus}{}

\noindent\textit{Papers, monographs, systems, software, visualizations,
and artifacts constituting the body of work.}

\bigskip

\noindent The corpus is organized below by type rather than by chronology
or program affiliation. Many works belong to multiple programs
simultaneously. Where that is the case, the primary program association
is indicated. Cross-program relationships are noted where they are
structurally important rather than merely thematic.

The corpus is not complete. It is a record of the work to the present
moment. New instruments are still being built. New habitats are still
being constructed. New papers are accumulating. The list below should be
read as a survey of the terrain as it currently stands, not as a
closed inventory.

\section{Major Monographs}

\subsection{Persistence Before Truth} The foundational monograph of the
Ontological Critique Program. Argues that persistence is a more primitive
category than truth: a structure must persist long enough to be evaluated
before the question of its truth or falsity arises. Develops the formal
framework of distinction preservation and introduces the distinction field
as a primary theoretical object.

\subsection{The Fate of Distinguishability} The second volume of the
Distinction Trilogy. Investigates the conditions under which distinguishability
is preserved across transformations, compressed representations, noisy
channels, and projection operations. Introduces the collapse quotient and
develops the geometry of distinction boundaries.

\subsection{The Ecology of Distinctions} The third volume of the Distinction
Trilogy. A textbook developing the biological and
ecological metaphors of distinction theory. Covers distinction niches,
distinction food webs, extinction pressures on conceptual boundaries,
and the conditions under which new distinctions emerge and stabilize.
Includes a fourteen-term index and extended worked examples drawn from
the Pullman company town, the Denver Post, and the discovery of
\textit{H. pylori}.

\subsection{Constraint, Projection, and Reachability} The CPR monograph.
Ninety-two chapters, four hundred and thirty-five pages. The most
comprehensive single treatment of the admissibility framework. Develops
the full formal apparatus of admissibility manifolds, reachability volumes,
projection operators, and constraint propagation. The primary technical
reference for the Admissibility Program.

\subsection{MEM\textbar8} A monograph on memory as recoverable continuation.
Develops the history-centered computational architecture and the formal
account of ecphory as reconstruction rather than retrieval. Argues that
memory is a geometric property of event histories and that the conventional
storage model is an artifact of technological rather than cognitive
constraints.

\subsection{Holonomic Space} A one-hundred-and-ninety-five-page textbook
on the geometry of path-dependent systems. Develops holonomy as the primary
mathematical tool for systems in which the current state is insufficient to
determine future behavior without reference to the history of the trajectory.

\subsection{The Flyxion Atlas} A meta-monograph cataloguing the full
Flyxion corpus. At the time of the present edition: approximately two
thousand one hundred and fifty lines, sixty-five entries (figures reflect
a snapshot; the Atlas is a living document). Organized by framework,
program, and instrument type.
The present document is a successor to and expansion of the Atlas.

\section{Repair-Theoretic Series}

Five papers developing the formal foundations of Repair Theory.

\subsection{Observability Is Restorability} Argues that the correct account
of observability is not epistemic but repair-theoretic. A state is
observable to the degree that it is restorable from available evidence.
Develops the formal connection between observation and reconstruction.

\subsection{Erasure Is Deformation} Argues that information erasure is not
annihilation but geometric deformation of a distinction structure. Erased
information leaves traces in the shape of the space that remains. Develops
the formal geometry of erasure marks.

\subsection{Admissibility Before Truth} Companion paper to Persistence
Before Truth. Argues that the admissibility of a proposition — whether it
can function within the current constraint structure — is prior to its
truth value. A proposition that cannot be evaluated within the current
admissibility structure is not false; it is inadmissible.

\subsection{The Geometry of Witnesses} Develops the formal theory of
witnesses as geometric objects. A witness is not merely evidence. It is
a structure that preserves the shape of a distinction boundary against
the pressure to collapse it. Develops witness algebras and witness
persistence measures.

\subsection{Conservation of Ambiguity} Argues that ambiguity is conserved
across transformations in a formally precise sense. When a transformation
removes ambiguity from one region of a distinction structure, it introduces
equivalent ambiguity elsewhere. Develops the formal analogue of conservation
laws for distinction boundaries.

\section{RSVP Papers and Monograph Chapters}

\subsection{Gravity as Entropy Descent} Develops the RSVP account of
gravity as a consequence of scalar entropy field gradients rather than
spacetime curvature. Derives the Newtonian limit and discusses deviations
from general relativity at cosmological scales.

\subsection{RSVP Cosmology} A monograph-in-progress developing the full
RSVP framework. Covers CTOV equations, scalar entropy fields, vector baryon
flows, structure formation, cosmic microwave background reintegration,
alternatives to inflation, and experimental predictions.

\subsection{Repairing Futures} A paper connecting the RSVP framework to
the PHYSIFORMER neural architecture. Argues that the PHYSIFORMER's
approach to physical simulation instantiates admissibility-theoretic
principles at the computational level. Develops the formal correspondence
between entropic smoothing in RSVP and repair operations in the Repair
Theory Program.

\subsection{Emergent Structures and Control in Neural and Cosmic Systems}
A comparative analysis of emergence in neural networks and cosmological
structure formation. Argues that both domains exhibit the same signature:
local repair operations producing global coordination without centralized
control.

\section{Negation and Inferential Field Theory}

\subsection{Negation Before Logic} A monograph arguing that negation is
not a logical operation defined within a formal system but an orientation
transformation defined within an inferential field. A negation does not
negate a proposition. It reorients the inferential trajectory of a system
relative to the proposition. Develops inferential field theory as an
alternative to propositional logic for modeling reasoning under constraint.

\subsection{Negation as Orientation Transformation} The formal monograph
accompanying Negation Before Logic. Develops the full mathematical
framework of inferential fields, orientation operators, and trajectory
reorientation. The primary technical reference for the negation program.

\section{Cognition and AI Papers}

\subsection{Against Latent Fundamentalism} A critique of the assumption
that latent representations are the fundamental objects of machine learning.
Argues that latent spaces are derivative of the distinction structures that
generate them and that treating latent representations as primary obscures
the admissibility structure of the learning process.

\subsection{Procedural Generation of Attention} A paper in the intersection
of film theory and cognitive science. Argues that cinema functions as a
system for procedurally generating viewer gaze trajectories through
admissibility structures. Develops the formal connection between narrative
geometry, attention routing, and the Procedural Attention Program.

\subsection{The Complex Numbers in Quantum Mechanics} An admissibility-theoretic
reading of the Barrios et al.\ result. Argues that the necessity of complex
numbers in quantum mechanics follows from an admissibility constraint on
the distinction structures available to a measuring apparatus, rather than
from an independent postulate about the mathematical structure of physical
law.

\section{Political Economy and Institutional Papers}

\subsection{Fiscal Reachability} Applies the admissibility framework to
public finance. Argues that the most important property of a fiscal system
is not its current budget position but its reachability volume: the set
of future fiscal states it can still access without structural rupture.
Develops formal measures of fiscal admissibility and applies them to
austerity, public investment, and institutional exhaustion.

\subsection{Platform Extraction and Distinction Collapse} Applies the
distinction framework to platform economics. Argues that platform
enshittification is formally a process of distinction collapse: the
distinctions between producer, consumer, and platform that originally
structured the ecosystem erode under extraction pressure, destroying
the admissibility structure that made the platform valuable.

\subsection{Civilizations as Hypotheses} An essay arguing that civilizations
are best understood not as entities but as sustained hypotheses about which
distinctions are worth maintaining. Civilizational collapse is the
withdrawal of that hypothesis.

\section{Internet Archaeology Papers}

\subsection{GitHub Follow-Network Survival Curves} An empirical study of
developer activity lifecycles in the GitHub social network. Analyzes
longitudinal follow-network data to identify survival curves for developer
engagement. Finds a U-shaped survival pattern: high early attrition,
stable middle-career persistence, and late-career decline. One of the
few works in the corpus that moves from theory to direct measurement.

\subsection{Bot-Farm Investigations} A series of analyses of coordinated
inauthentic behavior in open-source communities. Identifies the
distinguishing signatures of bot-farm activity in repository interaction
patterns and develops detection heuristics based on distinction geometry:
genuine users generate heterogeneous distinction patterns while coordinated
bots generate homogeneous ones.

\subsection{Repository Ecosystem Dynamics} An analysis of attention and
abandonment dynamics in open-source software ecosystems. Applies the
admissibility framework to software project lifecycles, identifying
the conditions under which a project loses its continuation capacity
and becomes effectively abandoned despite remaining technically accessible.

\section{Software and Systems}

\subsection{Spherepop interpreter} A working implementation of the
Spherepop irreversible event calculus. Written in C. Implements refusal
as a first-class operation and maintains full event histories as
runtime data structures.

\subsection{Icepick} A codebase navigation and snapshot system. Implements
traversal history preservation for software repositories. Under active
development.

\subsection{cprsh} A constraint-first shell environment. Implements the
CPR ontology at the level of interactive computation.

\subsection{Audioscope} A real-time audio visualization system with CRT
and phosphor aesthetics. Supports VTT subtitle overlay and Lissajous
trace rendering.

\subsection{Lava Thought Visualizer} A semantic terrain visualization
system. Renders conceptual proximity as navigable topographic landscape.

\subsection{Topological Loop Monitor} A real-time visualization system
for topological features in dynamical data streams.

\subsection{Blastoids} An evolutionary image generation system implementing
aesthetic selection over morphological search spaces.

\subsection{flyxion\_normalize.sh} An automation script for normalizing
transcription errors across large document corpora. Part of the document
ecology infrastructure.

\subsection{hotstrings.ahk} An AutoHotKey script implementing a custom
hotstring system for accelerated input of specialized notation, including
Arabic script, Phoenician characters, and symbolic writing systems.

\section{Creative Works}

\subsection{City of Brutes} A screenplay. One hundred and fifteen pages.
Features the Library of Revisions and the Library of Contradictions as
primary institutional settings. See Section V for methodological discussion.

\subsection{Notes from the Collapsewell} An ergodic novel. One hundred
and two pages. Implements intentional spelling degradation across the
course of the text. See Section V for methodological discussion.

\subsection{Worlds of If} A speculative documentary magazine series.
Multiple issues. Each investigates a reachable future that was not reached.

\subsection{The Flyxion Instrumentarium} A bestiary of conceptual entities.
Ongoing. See Section V for methodological discussion.

\subsection{Nine Parsecs} A graphic novel in development. Long-form
sequential art project.

\subsection{Flower Wars} A screenplay. Historical speculative fiction.

\subsection{The Incoherence} A screenplay. Philosophical drama.

\subsection{Yarncrawlers of Titan} A comic series. Science fiction.
Developed in parallel with the Yarncrawler framework paper.

\subsection{The Flyxion Atlas Cards} A fifty-two-card educational set.
Each card presents a single concept, instrument, or program entry from
the corpus in a format designed for transmission across technological
and institutional contexts.

\subsection{Flyxion Instrumentarium: Medieval Bestiary Project} Extended
bestiary entries in the style of medieval natural history manuscripts.
Occupies the intersection of historical document aesthetics and
contemporary theoretical content.

\section{Reference and Infrastructure}

\subsection{The Flyxion Atlas} The primary corpus catalog. Approximately
sixty-five entries across all program areas (snapshot figures; see note
above). The present atlas supersedes
and extends it.

\subsection{Spherepop repository catalog} A catalog of interconnected
textbooks and essays hosted at standardgalactic.github.io/spherepop.
Each entry carries a one-to-two sentence description situating the work
within the broader program structure.

\subsection{GitHub repository ecosystem} The primary distribution
infrastructure for software, papers, and related materials. Hosted
under the standardgalactic organization.

\bigskip\bigskip

\noindent\textit{The corpus described above is the evidence. The fifteen programs described in Section III are hypotheses about what the evidence means. The instruments described in Section IV are the devices through which evidence was collected. The habitats described in Section V are the environments in which the phenomena were allowed to occur.}

\medskip

\noindent\textit{Together they constitute not a completed investigation but a sustained one. The question that opened this atlas remains open. The work continues.}

\cleardoublepage
\appendix

% ---------------------------------------------------------------
% APPENDIX A — CHRONOLOGICAL TIMELINE
% ---------------------------------------------------------------
\chapter[Appendix A: Chronological Timeline]{Appendix A \\ Chronological Timeline}
\markboth{appendix a: chronological timeline}{}

\noindent\textit{A provisional record of when the major phases and projects
emerged. Dates are approximate. The timeline reflects the development of
the research as reconstructed from the corpus; it is not a publication
record.}

\bigskip

\section{Phase I — Representation (early period)}

The initial period centers on questions of information, representation,
and computation. Work in this phase investigates what information is
present in a system, how states are encoded, how machines process symbols,
and how meaning is represented in language and neural networks. The RSVP
framework begins taking shape as a physical substrate. Early drafts of
what will become the Spherepop calculus appear.

Key concerns: encoding, semantics, machine learning, computational state,
early cosmological speculation.

\section{Phase II — Distinction}

The vocabulary shifts toward distinction as a primitive. The Ontological
Critique Program opens. Early versions of distinguishability geometry
are developed. The claim that entities are derivative of distinctions
rather than distinctions being derivative of entities is first formalized.

The Ecology of Distinctions begins as a textbook project. The Distinction
Trilogy is conceived. RSVP is refined as a scalar-vector field theory.

Key works: early drafts of Persistence Before Truth; early Ecology of
Distinctions chapters; RSVP scalar field papers.

\section{Phase III — Repair}

The repair-theoretic turn. The five-paper repair series is developed:
Observability Is Restorability, Erasure Is Deformation, Admissibility
Before Truth, The Geometry of Witnesses, Conservation of Ambiguity.

Memory is reconceived as a repair process. Ecphory replaces retrieval
as the primary memory operation. MEM\textbar8 is developed as a
history-centered architecture.

The Holonomic Space textbook is completed (195 pages). City of Brutes
begins as a screenplay project.

Key works: five repair-theoretic papers; MEM\textbar8 monograph; Holonomic
Space; early City of Brutes drafts.

\section{Phase IV — Continuation}

Identity is relocated from substance to trajectory. The Continuation
Program formalizes path preservation and dynamic identity. Notes from
the Collapsewell begins as an ergodic novel experiment.

The Instrumentarium bestiary project opens. The four-vocabulary
civilization narrative framework (Inhabitants, Historians, Repairers,
Narrators) is developed. Worlds of If begins as a speculative documentary
series.

Spherepop reaches implementation as a working C interpreter. cprsh is
developed as a constraint-first shell.

Key works: Continuation Program papers; early Collapsewell drafts; Worlds
of If issues; Spherepop interpreter; Instrumentarium entries.

\section{Phase V — Admissibility}

The Admissibility Program becomes the dominant organizational framework.
The CPR monograph (Constraint, Projection, and Reachability) reaches
ninety-two chapters and four hundred and thirty-five pages. Admissibility
manifolds, reachability volumes, and the admissibility distortion metric
are formally developed.

RSVP reaches its mature form with CTOV equations. Gravity as Entropy
Descent is written. Repairing Futures connects RSVP to PHYSIFORMER.

The GitHub follow-network survival curve study is conducted — one of the
few empirical measurement projects in the corpus. Bot-farm investigations
follow.

Negation Before Logic and Negation as Orientation Transformation are
written. Procedural Generation of Attention is developed at the
intersection of film theory and cognitive science.

The Flyxion Atlas is assembled as the first corpus-level catalog
(approximately 2,153 lines, 65 entries).

Key works: CPR monograph; RSVP cosmology chapters; Gravity as Entropy
Descent; Repairing Futures; GitHub survival curve paper; Negation
monographs; Procedural Generation of Attention; Flyxion Atlas v1.

\section{Phase VI — Atlas Period (present)}

The center of gravity shifts toward synthesis. The survival question
crystallizes explicitly: \textit{how do structures survive the loss of
the conditions that originally produced them?} The question is identified
as having emerged independently from two directions — historical
reconstruction of the projects and classification of the programs — and
the convergence is treated as evidence of the question's depth.

The present document is the primary product of this phase. It supersedes
the earlier Flyxion Atlas and extends it into a full intellectual atlas
with program registry, instrument catalog, dependency graph, experimental
habitat documentation, corpus survey, and chronological record.

Ongoing: the Distinction Trilogy middle volume (The Fate of
Distinguishability); the Negation monograph series; new RSVP experimental
predictions; continued Collapsewell and Instrumentarium development;
Icepick system development; audio visualizer extensions.

\bigskip

\noindent\textit{This timeline will require revision as the corpus
continues to develop. It is offered as a snapshot of the sequence in
which the major conceptual moves occurred, not as a definitive account.
The phases overlap and the order is approximate. What the timeline records
is not dates but conceptual dependencies — which moves had to happen
before which other moves became possible.}

\cleardoublepage

% ---------------------------------------------------------------
% SECTION VII — ADMISSIBILITY: A SYNTHESIS
% ---------------------------------------------------------------
\chapter[Admissibility: A Synthesis]{Section VII \\ Admissibility}
\markboth{section vii: admissibility}{}

\noindent\textit{What the programs, instruments, and habitats reveal.}

\bigskip

\noindent The preceding sections have described a distributed investigation:
fifteen programs, dozens of instruments, numerous experimental habitats, and
a sprawling corpus. The risk is that this appears as an accumulation — a
collection of independent projects grouped only by their shared origin.

That appearance would be misleading.

The programs, instruments, and habitats converge on a single concept that
organized them long before any of them understood what was happening. That
concept is admissibility.

\section{The Definition}

Admissibility is the property of being compatible with the continuation of
a system.

An admissible state is a state from which some relevant future remains
accessible. An admissible proposition is a proposition that can be
entertained without destroying the distinction structure that makes
evaluation possible. An admissible transformation is a transformation that
preserves at least one pathway to a future the system values.

Admissibility is not truth. Truth is a property of propositions relative to
a world. Admissibility is a property of states, propositions, and
transformations relative to a trajectory. The trajectory is what matters.

Admissibility is not optimization. A system can be operating at peak
efficiency while its admissibility declines. Many historical collapses
occurred in societies that were locally optimal. The optimization obscured
the admissibility contraction until the contraction became irreversible.

Admissibility is not stability. A stable system can be fragile. Stability
over short time horizons is compatible with collapse over long ones. The
ancient Maya were stable until they were not. The collapse was not
precipitated by instability. It was precipitated by the exhaustion of
admissibility under accumulated stress.

Admissibility is the most fundamental property of any persisting system.
It is prior to truth, prior to optimization, prior to stability, prior to
all the properties conventional theory treats as primary.

\section{The Geometry}

Admissibility possesses a geometric structure.

Reachable futures form a volume in the state space of a system. That
volume has a shape, a boundary, an interior, an exterior. The boundary
separates admissible from inadmissible futures. The shape determines what
remains possible.

Admissibility manifolds are the mathematical objects that encode this
geometry. They are not fixed. They evolve with the system. Every decision
distorts the manifold. Every repair operation reopens a region of the
space. Every collapse deletes one.

The geometry of admissibility is the geometry of persistence. To understand
how a system continues — or fails to continue — one must understand how
its admissibility manifold evolves under the pressures it experiences.

\section{The Sources}

The admissibility framework was not invented. It was discovered in the
convergence of independent lines of inquiry.

The Ontological Critique Program discovered it by asking how distinctions
survive. Distinctions survive when they are admissible — when they can be
maintained within the current constraint structure without collapse.

The Repair Theory Program discovered it by asking how systems recover.
Recovery is the process of restoring admissibility to regions of the state
space that had become inaccessible.

The Memory Program discovered it by asking what memory preserves. Memory
preserves not representations but admissibility. A system that remembers
can reach futures that a system that has forgotten cannot.

The Computation-as-History Program discovered it by asking what the history
of a computation conveys. The history conveys the trajectory through
admissibility space — the record of which futures were accessed, which were
foreclosed, and which remain accessible.

The Admissibility Program discovered it by asking directly what it means
for a future to remain possible. The answer was geometry.

Even RSVP discovered it. The scalar entropy fields and vector baryon flows
are descriptions of admissibility evolution at cosmological scale. The
universe itself is a system whose admissibility manifold is the space of
possible cosmic futures. The CTOV equations govern how that manifold evolves.

The convergence was not planned. It was the phenomenon becoming visible
through multiple instruments at once.

\section{The Unity}

The fifteen programs are therefore not fifteen independent projects. They
are fifteen views of the same object.

The Ontological Critique Program studies how distinctions create
admissibility structures. The Admissibility Program studies the geometry
of those structures directly. The Repair Theory Program studies how they
are reconstructed after damage. The Continuation Program studies what it
means to move through them over time. The Memory Program studies how their
past shapes their present. The Coordination Geometry Program studies how
multiple agents negotiate shared admissibility structures. The Preference
Fields Program studies how values are defined over them. The Computation-as-History
Program studies how computations traverse them. The Procedural Attention
Program studies how experiences guide traversal. The Fictional Habitat
Program studies what it is like to be inside them. The Visualization
Program studies how they can be made visible. The Civilization and Futures
Program studies them at the largest scale. The Generative Evolution Program
studies how they can be expanded into regions previously uninhabitable.
RSVP studies them at the fundamental physical level.

Each program provides a different kind of access. The object is the same.

\section{The Practical Consequence}

The practical consequence of this convergence is that the work can now
be directed with a precision that was not previously possible.

Earlier work was exploratory. It followed the constraints of each local
inquiry: the demands of a given problem, the affordances of a given
instrument, the trajectory of a given investigation. That was necessary.
Without the exploration, the convergence could not have occurred.

But exploration has a cost. It scatters effort. It produces results that
are difficult to connect. It makes the shape of the overall enterprise
invisible to those inside it.

The recognition of admissibility as the organizing concept changes this.
It provides a principle of selection: future work can ask whether a project
contributes to the understanding of admissibility and, if so, at which
scale and through which instrument. It provides a principle of connection:
the relation between a work on fiscal policy and a work on the geometry of
latent spaces is not merely thematic — both contribute to the understanding
of admissibility in different domains, and the connection is structural.
It provides a principle of evaluation: a project is successful to the
degree that it makes admissibility more legible.

The investigation is no longer just a collection of work. It has become
a program.

\section{The Open Question}

Admissibility itself remains open.

The geometry is understood in outline but not in detail. The algebraic
properties of admissibility manifolds have not been fully characterized.
The dynamics of admissibility evolution are known only in special cases.
The relationship between local and global admissibility is poorly
understood. The connection between admissibility and preference — between
what is possible and what is valued — is the subject of ongoing work.

What is known is that admissibility is the right level of description.
The failures described in Section II all occurred because the wrong level
was chosen. Representation, prediction, optimization, state description,
information — all failed to capture what mattered. Admissibility did not
fail.

It remains to be seen how far the framework can be pushed. Whether it
can provide a unified account of persistence across all domains. Whether
the mathematics can be made rigorous. Whether the experimental habitats
can validate the formal claims. Whether the instruments can be refined
to the point of practical utility.

These questions are not at the end of the work. They are at the beginning.

\cleardoublepage

% ---------------------------------------------------------------
% SECTION VIII — CONTINUATION
% ---------------------------------------------------------------
\chapter[Continuation]{Section VIII \\ Continuation}
\markboth{section viii: continuation}{}

\noindent\textit{The trajectory of the work.}

\bigskip

\noindent The atlas closes where it began: with continuation.

The opening statement observed that every structure inherits a problem.
The conditions that produced it eventually disappear. The question is how
structures persist when those conditions are no longer available.

This is not a question to be answered once. It is a question to be lived
with. Each generation encounters it anew. The instruments change. The
habitats change. The formal frameworks change. The question remains.

The work described in this atlas is one response to that question. It is
not the only response. It is not the final response. It is the response
that emerged from a particular set of investigations, through a particular
sequence of failures and corrections, using a particular set of instruments
and habitats and formal frameworks.

The work continues.

\section{What Has Been Done}

Fifteen programs have been initiated. The full apparatus of admissibility
geometry has been developed in outline. A comprehensive instrument catalog
has been assembled. Experimental habitats have been constructed that allow
the phenomenon to be directly inspected. The corpus has grown to include
monographs, papers, software, visualizations, and creative works.

The convergence has been identified. The organizing concept has been named.
The failures that preceded the convergence have been documented.

This is not nothing. It is substantial. It is enough to orient future work
and to connect the scattered inquiries that preceded it.

\section{What Remains}

What remains is everything that follows from the recognition.

The formal framework must be developed in detail. The algebraic structure
of admissibility manifolds must be characterized. The dynamics must be
understood. The connection to preference must be established. The
experimental habitats must be refined and their results integrated into
the theory. The instruments must be improved. The corpus must continue
to accumulate.

But the work is now oriented. The exploratory phase is over. The
programmatic phase has begun.

\section{The Completion}

The atlas is complete. The work is not.

\medskip

\noindent The question remains open. That is why the work continues.

\bigskip

\begin{center}
\textsc{Flyxion}

\medskip

\textit{The trajectory}
\end{center}

\cleardoublepage

% ---------------------------------------------------------------
% BACK MATTER
% ---------------------------------------------------------------
\pagestyle{plain}

\chapter*{Colophon}
\addcontentsline{toc}{chapter}{Colophon}
\markboth{colophon}{}

\noindent This volume was composed in \textit{TeX Gyre Pagella}, a
reinterpretation of Hermann Zapf's Palatino. The typeface was chosen for
its readability across future reading environments and its
survivability — qualities aligned with the subject of the work itself.

\medskip

\noindent The text was set in a 6\thinspace$\times$\thinspace9 inch trim
with generous margins to accommodate annotation, marginalia, and the kind
of active reading that the content demands. The layout follows the
principles of readability-first publishing: legibility over novelty,
survival over presentation.

\medskip

\noindent The document is designed to be archived in multiple formats:
the present PDF, source code in a version-controlled repository,
plain-text extraction, and \LaTeX{} source that can be recompiled as
tools and standards evolve.

\medskip

\noindent The colophon itself is a persistence structure: the part of the
document that tells future readers how to understand the material object
they hold and how to preserve it.

\medskip

\noindent\textit{The work continues.}

\cleardoublepage

\chapter*{Index}
\addcontentsline{toc}{chapter}{Index}
\markboth{index}{}

\noindent\textit{Page numbers are approximate and reflect the state of the
atlas at the time of compilation. The index will be revised in subsequent
editions as the corpus expands.}

\bigskip

\noindent\textsc{Admissibility} (geometry, manifolds, volumes, distortion,
equivalence), passim

\medskip
\noindent\textsc{Admissibility Program}, 5, 7, 26, 33, 44

\medskip
\noindent\textsc{Audioscope}, 7, 12, 14, 24

\medskip
\noindent\textsc{Blastoids}, 8, 12, 14, 25

\medskip
\noindent\textsc{City of Brutes}, 7, 16, 18, 25

\medskip
\noindent\textsc{Civilization and Futures Program}, 7, 27, 44

\medskip
\noindent\textsc{CLIO}, 11, 14

\medskip
\noindent\textsc{Continuation Program}, 6, 11, 17, 44

\medskip
\noindent\textsc{Coordination Geometry Program}, 6, 44

\medskip
\noindent\textsc{CPR} (Constraint, Projection, Reachability), 20, 24

\medskip
\noindent\textsc{cprsh}, 11, 14, 24

\medskip
\noindent\textsc{Distinction Ecology}, 17

\medskip
\noindent\textsc{Distinction Trilogy}, 19, 20

\medskip
\noindent\textsc{Fictional Habitat Program}, 7, 27, 44

\medskip
\noindent\textsc{Flyxion Atlas}, 20, 25, 26

\medskip
\noindent\textsc{Flyxion Instrumentarium}, 16, 17, 25

\medskip
\noindent\textsc{Generative Evolution Program}, 7, 8, 44

\medskip
\noindent\textsc{History machines}, 11, 14

\medskip
\noindent\textsc{HYDRA}, 11, 14

\medskip
\noindent\textsc{Icepick}, 6, 11, 14, 24

\medskip
\noindent\textsc{Lava Thought Visualizer}, 7, 12, 14, 24

\medskip
\noindent\textsc{Marqants}, 6, 11, 14

\medskip
\noindent\textsc{Memory Program}, 6, 11, 17, 44

\medskip
\noindent\textsc{MEM\textbar8}, 11, 14, 20

\medskip
\noindent\textsc{Notes from the Collapsewell}, 7, 16, 17, 25

\medskip
\noindent\textsc{Ontological Critique Program}, 5, 17, 19, 20, 43, 44

\medskip
\noindent\textsc{Preference Fields Program}, 6, 44

\medskip
\noindent\textsc{Procedural Attention Program}, 7, 12, 22, 44

\medskip
\noindent\textsc{Repair Theory Program}, 5, 6, 10, 20, 21, 22, 43, 44

\medskip
\noindent\textsc{RSVP} (Relativistic Scalar-Vector Plenum), 8, 9, 10, 14, 21, 43, 44

\medskip
\noindent\textsc{Spherepop}, 6, 11, 14, 24, 26

\medskip
\noindent\textsc{TARTAN}, 11, 14

\medskip
\noindent\textsc{Topological Loop Monitor}, 7, 12, 14, 24

\medskip
\noindent\textsc{Visualization Program}, 7, 27, 44

\medskip
\noindent\textsc{Worlds of If}, 7, 16, 18, 25

\medskip
\noindent\textit{The work continues.}

\end{document}
